\chapter{Implementación del algoritmo}

% Capítulo descriptivo: dónde encaja el algoritmo en el flujo composicional,
% cómo se construyó, qué decisiones de diseño se tomaron por sobre lo que
% deja abierto el paper, qué adaptaciones requirió el contexto de síntesis
% de controladores, y cómo se validó la correctitud.


\section{Introduccion}
% Una frase sobre qué es MTSA y por qué es la plataforma natural donde
% implementar el reemplazo de WSOE por branching bisimilarity.
La herramienta MTSA \cite{mtsa} es el framework de desarrollo realizado en Java para varias iniciativas de LaFHIS en el contexto de la investigacion sobre LTSs. Es por eso que tanto el algoritmo composicional presentado por Uchitel y Gagliardi \#AGREGAR CITA\# como la version de WSOE utilizada para la composicionalidad fueron implementados en esta herramienta. Era entonces natural realizar el desarrollo del algoritmo de Branching Equivalence sobre el mismo framework, ante la promesa de que esto reduciria la complejidad. A su vez, nos permitira comparar de forma mas veridica ambos algoritmos, realizando pruebas de Benchmarks tanto para tiempos de ejecucion como para el uso de memoria.

En primer lugar, se realizo la implementación del algoritmo de Branching Equivalence en el archivo BranchingEquivalence.java, creado en mtstools/DCS/Compositional/ ya que bajo este modulo estan las herramientas utilizadas por el algoritmo composicional.


A su vez, en CompositionalApproach.java se reemplazó la llamada original al algoritmo de WSOE por la llamada a la nueva funcion que utiliza el algoritmo de Groote y Jansen, reemplazando tambien los preprocesamientos necesarios para los algoritmos que se hacen en algunos casos fuera de las funciones principales.

La llamada al algoritmo de Branching Equivalence queda entonces de esta forma
\begin{lstlisting}[language=Java, numbers=none]                                
  Pair<List<Set<Long>>, List<Set<Triple<Long, String, Long>>>> branchingPartition
                    = BranchingEquivalence.getPartitions(localControllerMTS, localAlphabet, totalTranslator, goal.getFluents());

  MTS<Long, String> minimizationResult = BranchingEquivalence.buildMinimisedMTSFromPartition(localControllerMTS, tauLabels, translatorControllable, branchingPartition);
\end{lstlisting}                                                               

donde en la funcion getPartitions(...) se implementa el algoritmo de Groote y Jansen y se devuelve la partición final de estados y transiciones, y en buildMinimisedMTSFromPartition(...) se construye el MTS minimizado a partir de la partición calculada.

A su vez, el algoritmo recibe como parámetros el MTS a minimizar, el alfabeto local, el traductor total y los fluents objetivo. 

En primer lugar, la modalidad de desarrollo fue en primera mano implementar la logica del algoritmo lo mas fielmente al paper posible, es decir, sin tomar en cuenta las optimizaciones ni las cotas de complejidad. En la version cero (v0), el algoritmo 

\section{Primer desarrollo del algoritmo (v0)}


\subsection{Idea general: dos particiones que se refinan en simultáneo}
Como fue explicado en el capitulo 3, el algoritmo de Groote y Jansen mantiene dos estructuras de refinamiento: una partición de estados $\Pi_s$ y una particion de transiciones $\Pi_t$. En $\Pi_s$ se encuentran los "bloques" de estados estables tal que si dos estados estan en distintos bloques entonces NO son branching equivalent. En $\Pi_t$ se encuentran los "bunches" de transiciones no-inertes, es decir, las transiciones que pueden distinguir estados. Ambas estructuras se refinan a la par: cada vez que un bloque de estados se parte en dos, se actualizan los bunches para reflejar el nuevo bloque, y cada vez que un bunch se parte en slices por (acción, bloque destino), se generan nuevos splitters para estabilizar los bloques afectados. El proceso continúa hasta llegar al punto fijo donde no quedan ni splitters pendientes ni bunches por refinar: las clases de equivalencia son entonces los bloques finales de $\Pi_s$.
La idea para esta primera implementación fue mantener la estructura del paper, donde el loop principal itera sobre cada conjunto de trancisiones no triviales que estan en $\Pi_t$, tomando un subconjunto chico de transiciones. Este subconjunto de trancisiones, es separado de su "bunch" original y añadido a $\Pi_t$ como un nuevo bunch. Luego, se deben refinar los bloques de $\Pi_s$ para que queden estables respecto al nuevo conjunto de trancisiones obtenido.

\subsection{Estructuras utilizadas}
$\Pi_s$ se implementó como \lstinline|List<Set<Long>>|, es decir una lista de conjuntos de estados. Cada \lstinline|Set<Long>| representa un bloque de estados.
$\Pi_t$ se implementó como \lstinline|List<Set<Triple<Long, String, Long>>>|, es decir una lista de conjuntos de transiciones. Cada \lstinline|Set<Triple<Long, String, Long>>| representa un bunch de transiciones no-inertes.
A su vez, se debió utilizar una estructura \lstinline|Pi_t_cola| para poder iterar sobre los bunches a la vez que se actualizaba la lista de pendientes a refinar.
Tal como se menciona en el paper de Groote y Jansen, se utiliza tambien la estructura del \lstinline|splitterList| para almacenar los splitters generados durante el proceso de refinamiento. A su vez, se creo la clase \lstinline|Splitter| para representar cada splitter, que contiene el bloque que se va a partir, el conjunto de transiciones que cruzan bloques y un flag de primary/secondary, un conjunto de transiciones marcadas para identificar que esos bloques son estables y un groupId para identificar que Splitters se originaron a partir de la misma division inicial.

\subsection{Inicialización}
En primer lugar, el primer preprocesamiento requerido para este algoritmo es eliminar las componentes tau fuertemente conexas. La implementacion de este paso usa la logica del algoritmo de Kosaraju, que tiene una complejidad O(n + m), por lo que cumple con la complejidad prometida por Groote y Jansen.
Luego, se debe inicializar $\Pi_s$, que arranca con dos bloques (Bvis y Binvis), donde Bvis contiene el conjunto de estados desde los cuales se puede alcanzar una transición visible y Binvis es su complemento.
Para realizar esta inicializacion, se ejecuta la funcion \texttt{computeBvis} que genera un grafo de predecesores, almacenando para cada transicion, cual fue su origen. Luego, identifica cuales son visibles y realiza un DFS para propagar las acciones visibles y identificar desde que estados son alcanzables. Por lo tanto, la complejidad resultante de la funcion es tambien O(n + m). Finalmente, Bvis se inicializa con el conjunto de estados identificados por esta funcion como visibles y Binvis se inicializa con el resto.
$\Pi_t$ es inicializado con un único bunch que contiene todas las transiciones no-inertes.


\section{Optimizaciones y mejoras (v1)}
\subsection{Estructura del bucle: dos fases que alternan}
% El bucle principal alterna dos fases: Fase 1 estabiliza estados drenando
% splitterList; Fase 2 refina un bunch y genera nuevos splitters. Continúa
% hasta que ambas estructuras quedan vacías.


\subsection{Fase 1: estabilización de estados}
% Por cada splitter pendiente se ejecuta split(B, transitions, marks); si
% parte el bloque B en R y U se actualizan las estructuras, se subdividen los
% splitters obsoletos que apuntaban a B y se propaga la cascada de $\tau$
% no-inertes nuevas mediante la cola newFrontiers.


\subsection{Fase 2: refinamiento de bunches}
% Se elige un bunch T de la cola, se lo parte en slices por (acción, bloque
% destino), se elige una small half de tamaño $\le |T|/2$ y se generan
% splitters primary/secondary asociados por groupId para que la próxima
% Fase 1 los procese.


\subsection{Construcción del MTS minimizado}
% Cada bloque recibe un identificador nuevo, cada transición $s \xrightarrow{a} s'$
% del MTS original se mapea a $\text{block}(s) \xrightarrow{a} \text{block}(s')$
% en el resultado, y los self-loops $\tau$ controlables se renombran a
% c_<acción> y se registran en translatorControllable.


\section{Decisiones de diseño no dictadas por el paper}

\subsection{Dos fases en lugar del nesting outer-while + inner-for}
% El paper anida un while externo sobre bunches con un for interno sobre
% splitters. Implementar dos fases que alternan al mismo nivel no cambia la
% complejidad asintótica, pero da invariantes verificables y separa los
% efectos de la restabilización del flujo principal.


\subsection{Identidad por referencia para la cola de bunches}
% Los bunches se mutan durante el algoritmo (addAll, removeAll). Una cola
% basada en equals/hashCode de Set produce falsos negativos en contains tras
% la mutación. Una cola que decide pertenencia por identidad de referencia
% (IdentityQueue) evita el bug y permite mutar bunches encolados con seguridad.


\subsection{Índice inverso estado destino $\to$ bunches}
% La estructura targetStateToBunches mapea cada estado destino al conjunto
% de bunches que lo contienen. Cuando se parte un bloque, esto permite
% re-encolar sólo los bunches afectados en lugar de re-escanear todo $\Pi_t$.


\subsection{Subdivisión sistemática de splitters obsoletos}
% Cuando un splitter parte el bloque B en R y U, los demás splitters
% pendientes que apuntaban a B quedan sin un blanco válido. La rutina
% refineSplitters los subdivide en dos splitters (uno para R, uno para U)
% preservando marcas según primary/secondary. El paper trata este caso de
% manera informal.


\subsection{SCCs como unidad atómica en split}
% R arranca con las SCCs enteras de los estados origen de las marcas, no con
% estados sueltos. Esto refleja una propiedad del algoritmo de Groote: los
% estados de una misma SCC $\tau$-conexa son siempre branching-equivalentes
% entre sí, por lo que separarlos en R y U sería incorrecto.


\subsection{Detección de bottom states por SCC}
% Una SCC se considera bottom si ninguno de sus estados tiene una transición
% $\tau$ saliente que termine en otra SCC del mismo bloque. Trabajar a nivel
% de SCC en lugar de estado a estado mejora tanto la corrección semántica
% como la eficiencia.


\subsection{Hashing de bloques en tiempo constante}
% El hashCode de Set<Long> es $O(n)$ y domina el tiempo de Fase 2 cuando los
% bloques son grandes. Asignar un identificador entero por bloque
% (blockIdMap) permite indexar las slices por Pair<String, Integer> y reduce
% el costo de hashing a $O(1)$.


\subsection{Restabilización en ambas direcciones}
% Tras un split en R y U, el paper sólo contempla nuevas $\tau$ no-inertes
% en sentido R \to U. En MTSs reales también aparecen en sentido U \to R.
% La cola newFrontiers propaga la cascada en ambas direcciones, agendando
% pares (R, U) y (U, R) hasta agotar las nuevas no-inertes generadas.


\section{Adaptaciones al contexto composicional}

\subsection{Estado de error como bloque inicial separado}
% El estado de error ($-1$) se trata como un bloque inicial propio en $\Pi_s$.
% Mezclarlo con Bvis o Binvis es semánticamente incorrecto: en el contexto
% de DCS el estado de error no es equivalente a ningún estado de operación
% normal.


\subsection{Self-loops $\tau$ controlables}
% Si un self-loop $\tau$ corresponde a una acción controlable, en el MTS
% minimizado se renombra a c\_<acción> y se registra en translatorControllable.
% Sin esta intervención la acción se perdería al colapsar el self-loop.


\subsection{Initiating actions de fluents fuera de $\tau$}
% Una acción que dispara un fluent modifica el valor del fluent y por lo
% tanto distingue estados desde la lógica del sistema. Tratarla como $\tau$
% haría que el algoritmo colapse estados que el modelo considera distintos,
% por lo que las initiating actions se sacan de tauLabels antes del bucle.


\subsection{Reuso de partición precomputada}
% La API buildMinimisedMTSFromPartition permite construir el MTS minimizado
% a partir de una partición ya calculada, sin recomputar la equivalencia.
% Es necesaria para que el flujo composicional pueda intercalar cómputos y
% reutilizar resultados.


\subsection{Eventos controlables, no-controlables y estados marcados}
% El algoritmo del paper opera sobre LTS "puros". Acá hay que preservar
% durante la minimización la información extra que aporta el contexto:
% qué eventos son controlables y cuáles no, y qué estados están marcados.


\section{Validación de correctitud}

\subsection{Estrategia}
% ¿Comparaste salidas con WSOE en casos chicos? ¿Verificaste alguna
% propiedad invariante (por ejemplo, que el LTS minimizado sea
% branching-bisimilar al original)? ¿Probaste sobre los casos del paper
% composicional?


\subsection{Tests automatizados}
% ¿Hay tests JUnit? ¿Cuántos casos? ¿Dónde viven en el repo?


\subsection{Casos de borde cubiertos}
% $\tau$-loops, MTSs sin transiciones invisibles, bloque error, fluents,
% bloques con un único estado, transiciones controlables que se vuelven
% self-loops tras la minimización.
